\begin{thebibliography}{99}\small
\bibitem{allen} C. S. Allen, Q. Chen, V. L. Willson, and J. N. Hughes, ``Quality of research design moderates effects of grade retention on achievement: A meta-analytic, multilevel analysis,'' \textit{Educational Evaluation and Policy Analysis}, vol. 31, no. 4, pp. 480--499, 2009. \href{https://doi.org/10.3102/0162373709352239}{doi:10.3102/0162373709352239}.
\bibitem{goos} M. Goos, J. Pipa, and F. Peixoto, ``Effectiveness of grade retention: A systematic review and meta-analysis,'' \textit{Educational Research Review}, vol. 34, art. 100401, 2021. \href{https://doi.org/10.1016/j.edurev.2021.100401}{doi:10.1016/j.edurev.2021.100401}.
\bibitem{tomchin} E. M. Tomchin and J. C. Impara, ``Unraveling teachers' beliefs about grade retention,'' \textit{American Educational Research Journal}, vol. 29, no. 1, pp. 199--223, 1992. \href{https://doi.org/10.3102/00028312029001199}{doi:10.3102/00028312029001199}.
\bibitem{bonvin} P. Bonvin, ``The role of teacher attitudes and judgement in decision-making: The case of grade retention,'' \textit{European Educational Research Journal}, vol. 2, no. 2, pp. 277--294, 2003. \href{https://doi.org/10.2304/eerj.2003.2.2.6}{doi:10.2304/eerj.2003.2.2.6}.
\bibitem{klapproth} F. Klapproth and C. Beyte, ``GPA, student gender, school subject, and teacher attitudes in grade retention decisions: An experimental study,'' \textit{European Journal of Psychology of Education}, vol. 41, art. 24, 2026. \href{https://doi.org/10.1007/s10212-026-01072-7}{doi:10.1007/s10212-026-01072-7}.
\bibitem{pajares} M. F. Pajares, ``Teachers' beliefs and educational research: Cleaning up a messy construct,'' \textit{Review of Educational Research}, vol. 62, no. 3, pp. 307--332, 1992. \href{https://doi.org/10.3102/00346543062003307}{doi:10.3102/00346543062003307}.
\bibitem{conjoint} J. Hainmueller, D. J. Hopkins, and T. Yamamoto, ``Causal inference in conjoint analysis: Understanding multidimensional choices via stated preference experiments,'' \textit{Political Analysis}, vol. 22, no. 1, pp. 1--30, 2014. \href{https://doi.org/10.1093/pan/mpt024}{doi:10.1093/pan/mpt024}.
\bibitem{validation} J. Hainmueller, D. Hangartner, and T. Yamamoto, ``Validating vignette and conjoint survey experiments against real-world behavior,'' \textit{Proceedings of the National Academy of Sciences}, vol. 112, no. 8, pp. 2395--2400, 2015. \href{https://doi.org/10.1073/pnas.1416587112}{doi:10.1073/pnas.1416587112}.
\bibitem{prediction} S. Mullainathan and J. Spiess, ``Machine learning: An applied econometric approach,'' \textit{Journal of Economic Perspectives}, vol. 31, no. 2, pp. 87--106, 2017. \href{https://doi.org/10.1257/jep.31.2.87}{doi:10.1257/jep.31.2.87}.
\bibitem{rf} L. Breiman, ``Random forests,'' \textit{Machine Learning}, vol. 45, pp. 5--32, 2001. \href{https://doi.org/10.1023/A:1010933404324}{doi:10.1023/A:1010933404324}.
\bibitem{shap} S. M. Lundberg and S.-I. Lee, ``A unified approach to interpreting model predictions,'' in \textit{Advances in Neural Information Processing Systems}, vol. 30, pp. 4765--4774, 2017. \url{https://papers.nips.cc/paper/7062-a-unified-approach-to-interpreting-model-predictions}.
\bibitem{nickell} S. Nickell, ``Biases in dynamic models with fixed effects,'' \textit{Econometrica}, vol. 49, no. 6, pp. 1417--1426, 1981. \href{https://doi.org/10.2307/1911408}{doi:10.2307/1911408}.
\bibitem{kunda} Z. Kunda, ``The case for motivated reasoning,'' \textit{Psychological Bulletin}, vol. 108, no. 3, pp. 480--498, 1990. \href{https://doi.org/10.1037/0033-2909.108.3.480}{doi:10.1037/0033-2909.108.3.480}.
\bibitem{spillane} J. P. Spillane, B. J. Reiser, and T. Reimer, ``Policy implementation and cognition: Reframing and refocusing implementation research,'' \textit{Review of Educational Research}, vol. 72, no. 3, pp. 387--431, 2002. \href{https://doi.org/10.3102/00346543072003387}{doi:10.3102/00346543072003387}.
\bibitem{post} J. M. Montgomery, B. Nyhan, and M. Torres, ``How conditioning on posttreatment variables can ruin your experiment and what to do about it,'' \textit{American Journal of Political Science}, vol. 62, no. 3, pp. 760--775, 2018. \href{https://doi.org/10.1111/ajps.12357}{doi:10.1111/ajps.12357}.
\bibitem{equivalence} D. Lakens, ``Equivalence tests: A practical primer for t tests, correlations, and meta-analyses,'' \textit{Social Psychological and Personality Science}, vol. 8, no. 4, pp. 355--362, 2017. \href{https://doi.org/10.1177/1948550617697177}{doi:10.1177/1948550617697177}.
\bibitem{holm} S. Holm, ``A simple sequentially rejective multiple test procedure,'' \textit{Scandinavian Journal of Statistics}, vol. 6, no. 2, pp. 65--70, 1979. \url{https://www.jstor.org/stable/4615733}.
\bibitem{lee} D. S. Lee, ``Training, wages, and sample selection: Estimating sharp bounds on treatment effects,'' \textit{Review of Economic Studies}, vol. 76, no. 3, pp. 1071--1102, 2009. \href{https://doi.org/10.1111/j.1467-937X.2009.00536.x}{doi:10.1111/j.1467-937X.2009.00536.x}.
\end{thebibliography}
